\chapter{Tools, Samples, and Migration}\label{app:tools}

\section{eiptool}\label{sec:eiptool}

\code{eiptool} is a command-line scanner and test adapter. It uses the host
port (POSIX or Windows). Its source is \file{examples/posix_eiptool/eiptool.c},
and it uses most of the scanner API.

\begin{table}[htbp]
\centering
\caption{eiptool commands}\label{tab:eiptool}
\small
\begin{tabularx}{\linewidth}{@{}P{0.50\linewidth}L@{}}
\toprule
\textbf{Command} & \textbf{Action} \\
\midrule
\code{eiptool scan} & Broadcast ListIdentity and list the devices. \\
\code{eiptool identity <ip>} & Read the identity of one device. \\
\code{eiptool tags <ip> [slot]} & List the tags of a controller. \\
\code{eiptool read <ip> <tag> [slot]} & Read a tag. \\
\code{eiptool write <ip> <tag> <type> <value> [slot]} & Write a tag. \code{type} is \code{dint}, \code{int}, \code{sint}, \code{real}, or \code{bool}. \\
\code{eiptool get <ip> <class> <instance> <attribute> [slot]} & Get\_Attribute\_Single. \\
\code{eiptool adapter} & Run the built-in test adapter until you stop it. \\
\bottomrule
\end{tabularx}
\end{table}

These rules apply:

\begin{itemize}
  \item \code{slot} is $-1$ by default.
  \item Numbers can be decimal or hexadecimal with \code{0x}.
  \item Set \code{EIP_NETIF=}\textit{name} to select the interface that the
        adapter reports and that the scan uses first.
  \item Explicit messages operate on \code{127.0.0.1}. Class~1 needs the
        address that the adapter reports.
\end{itemize}

\begin{lstlisting}[language=shell]
./build/eiptool adapter &
./build/eiptool get 127.0.0.1 4 101 3    # test assembly: 11 22 33 44 55 66 77 88
./build/eiptool get 192.168.1.10 0x01 1 7  # class instance attribute
./build/eiptool write 192.168.1.10 MyDint dint 42
\end{lstlisting}

\section{Examples in the Repository}

\begin{table}[htbp]
\centering
\caption{Example programs}\label{tab:examples}
\begin{tabularx}{\linewidth}{@{}P{0.34\linewidth}L@{}}
\toprule
\textbf{Directory} & \textbf{Shows} \\
\midrule
\file{examples/netburner_adapter} & An adapter on NetBurner, and an optional scanner tag read. \\
\file{examples/esp_idf_adapter} & An adapter as an ESP-IDF project. \\
\file{examples/clearcore_adapter} & A poll-mode adapter that controls ClearCore I/O. \\
\file{examples/posix_eiptool} & \code{eiptool}, with most of the scanner API. \\
\file{examples/eds} & A sample EDS file that agrees with the examples. \\
\file{examples/netburner_webui} & A browser test page. Application code, not part of the stack. \\
\bottomrule
\end{tabularx}
\end{table}

\section{Sample EDS File}\label{sec:edssample}

These are the device, assembly, and connection sections of
\file{examples/eds/QuickEIP_Adapter.eds}. They agree with the assemblies in
section~\ref{sec:start}. Vendor code 65000 is a placeholder.

\begin{lstlisting}[language=eds,basicstyle=\ttfamily\footnotesize]
[Device]
        VendCode = 65000;
        VendName = "QuickEIP Sample";
        ProdType = 43;
        ProdTypeStr = "Generic Device";
        ProdCode = 1;
        MajRev = 1;
        MinRev = 1;
        ProdName = "My Device";
        Catalog = "My Device";

[Assembly]
        Revision = 2;
        Object_Name = "Assembly Object";
        Object_Class_Code = 0x04;
        Number_Of_Static_Instances = 3;
        Assem100 =
                "Inputs",               $ device -> PLC (EIP_ASM_PRODUCED)
                "20 04 24 64 30 03",    $ path
                32,                     $ size in bytes
                0x0000,                 $ descriptor
                ,,                      $ reserved
                256,;                   $ member size in bits
        Assem150 =
                "Outputs",              $ PLC -> device (EIP_ASM_CONSUMED)
                "20 04 24 96 30 03",
                32,
                0x0001,
                ,,
                256,;
        Assem151 =
                "Config",               $ configuration, no data
                "20 04 24 97 30 03",
                0,
                0x0000,
                ,,
                0,Param2;

[Connection Manager]
        Connection1 =
                0x04010002,             $ class 1, cyclic, exclusive owner
                0x44640405,             $ O->T p2p, fixed, run/idle; T->O mcast or p2p
                Param1,32,Assem150,     $ O->T RPI, size, format
                Param1,32,Assem100,     $ T->O RPI, size, format
                ,,                      $ config part 1 (none)
                0,Param2,               $ config part 2: size, format
                "Exclusive Owner",      $ name
                "",                     $ help
                "20 04 24 97 2C 96 2C 64";  $ config 151, consume 150, produce 100

[Capacity]
        MaxIOConnections = 1;
        MaxMsgConnections = 2;
        TSpec1 = Rx, 32, 100;
        TSpec2 = Tx, 32, 100;
\end{lstlisting}

\section{Migration from QuickEIP 1.x}\label{sec:migration}

QuickEIP 2.0 replaces the NetBurner-only version 1.x with a portable library.
The adapter has these new functions: a configurable identity and
assemblies, a shared assembly data pool, multicast \TO, the TCP/IP Interface
and Ethernet Link objects, a settable IP configuration, Identity Reset, and
checks of the input. The sizes are now in \file{eip_config.h}.

\begin{procedure}{Change a 1.x NetBurner build to 2.0}
\begin{enumerate}
  \item Remove the old source lines for \file{eip_session.cpp},
        \file{eip_io.cpp}, and \file{eip_adapter_task.cpp} from the makefile.
  \item Add the source list of section~\ref{sec:netburner}.
  \item Change all timeouts from ticks to milliseconds (table~\ref{tab:migration}).
  \item Change \code{eip_adapter_start()} to \code{eip_adapter_start(0)}, or
        give a configuration.
  \item If you need the 1.x timeout, set \code{cfg.timeout_floor_ms = 2000}.
  \item If you use the web UI, add the files from
        \file{examples/netburner_webui/}.
\end{enumerate}
\end{procedure}

\begin{caution}
All timeouts are now in milliseconds. A value of \code{TICKS_PER_SECOND} (20)
that meant 1~s in version 1.x now means 20~ms. Examine each timeout argument.
\end{caution}

\begin{longtable}{@{}P{0.30\linewidth}P{0.30\linewidth}P{0.34\linewidth}@{}}
\caption{API changes from version 1.x}\label{tab:migration}\\
\toprule
\textbf{1.x} & \textbf{2.0} & \textbf{Notes} \\
\midrule
\endfirsthead
\toprule
\textbf{1.x} & \textbf{2.0} & \textbf{Notes} \\
\midrule
\endhead
\bottomrule
\endfoot
\code{EipOpenConfig.timeout_ticks} & \code{timeout_ms} & \code{TICKS_PER_SECOND * 5} becomes \code{5000}. \\
\code{eip_list_identity(ip, id, ticks)} & \code{(ip, id, ms)} & \\
\code{eip_scan_devices(d, n, ticks)} & \code{(d, n, ms)} & 0 still selects the default (3~s). \\
\code{eip_io_consume(io, d, n, ticks)} & \code{(io, d, n, ms)} & \\
\code{session.tcp_fd > 0} & \code{eip_session_is_open(\&session)} & The field is now \code{sock}, and 0 is a valid handle. \\
\code{io.udp_fd} & \code{io.sock} & \\
\code{EipSession.timeout_ticks} & \code{timeout_ms} & \\
\code{eip_adapter_start(void)} & \code{eip_adapter_start(const EipAdapterConfig *)} & Give 0 for the old behavior. \\
\code{EIP_MAX_PACKET} (enum) & Macro in \file{eip_config.h} & Same value. You can now change it. \\
\code{EIP_ADAPTER_ASSEMBLIES} (4) & \code{EIP_ADAPTER_MAX_ASSEMBLIES} (8) & The old name is still an alias. \\
\file{src/eip_port.h} & \file{src/eip_internal.h} & Internal header, new name. \\
\file{src/eip_adapter.h} & \file{include/eip_adapter.h} & Now public (the sans-I/O core). \\
\code{EipAdapterView} & More fields & Added \code{ip}, \code{multicast}, \code{rpi_us}, \code{timeout_ms}, \code{opens}, \code{timeouts}, \code{rx_dropped_*}, \code{rx_idle}, \code{udp_replies_dropped}, \code{tcp_clients}. \\
\end{longtable}

Changes of behavior:

\begin{itemize}
  \item The inactivity timeout now comes from the \OT{} RPI
        (equation~\ref{eq:timeout}). Version 1.x used the \TO{} interval and a
        2~s minimum.
  \item The Identity status now shows ``configured'' (table~\ref{tab:idstatus}).
  \item An unknown assembly instance answers CIP \hex{16}, not \hex{05}.
  \item Consumed data goes to all zeros on idle or timeout, as in version
        1.x. On idle, the adapter clears only the consumed point of the
        connection.
  \item The adapter accepts up to four TCP clients, not two. When all slots
        are busy, it drops a session that has been silent for 60~s.
  \item The 16-bit sequence count is still used for originators that send no
        Sequenced Address item. A repeated count keeps the connection open but
        does not replace the data. An older count is dropped.
\end{itemize}
